% !TEX root = ../main.tex
\section{Teoría elemental de números}

\subsection{Algoritmos de División y Euclides}

\begin{definition}[Divisibilidad]\label{def:divide}
    Sean $a, b \in \mathbb{Z}$ con $a \neq 0$. Decimos que $a$ \emph{divide} a $b$,
    y escribimos $a \mid b$, si existe $q \in \mathbb{Z}$ tal que $b = aq$.
    En ese caso también se dice que $a$ es un \emph{divisor} o \emph{factor} de $b$,
    o que $b$ es un \emph{múltiplo} de $a$.
    \src{1-1.1}
\end{definition}

\begin{proposition}[Propiedades de los enteros]\label{prop:ent}
    Para cualesquiera $a, b, c \in \mathbb{Z}$:
    \begin{enumerate}[label=(\arabic*), ref=\theproposition(\arabic*)]
        \item\label{prop:ent-cero} $0 \cdot a = 0$.
        \item\label{prop:ent-signo} $a(-b) = -(ab)$.
        \item\label{prop:ent-cancelacion} \emph{(Cancelación)} Si $a \neq 0$ y $ab = ac$,
              entonces $b = c$.
        \item\label{prop:div-multiplo} Si $a \neq 0$ y $a \mid b$, entonces $a$ divide
              a todo múltiplo de $b$: $a \mid bx$ para todo $x \in \mathbb{Z}$.
        \item\label{prop:div-transitiva} \emph{(Transitividad)} Si $a, b \neq 0$,
              $a \mid b$ y $b \mid c$, entonces $a \mid c$.
        \item\label{prop:div-comb-lineal} Si $a \neq 0$ divide a $b$ y a $c$, entonces
              divide a cualquier combinación lineal entera de ambos:
              $a \mid bx + cy$ para todo $x, y \in \mathbb{Z}$.
        \item\label{prop:div-orden} Si $a, b > 0$ y $a \mid b$, entonces $a \le b$.
        \item\label{prop:div-antisimetrica} Si $a, b \neq 0$, $a \mid b$ y $b \mid a$,
              entonces $a = \pm b$.
              \src{1-1.3}
    \end{enumerate}
\end{proposition}

\begin{theorem}[Algoritmo de la División para naturales]\label{teo:division}
    Sean $a$ y $b$ enteros donde $b > 0$. Entonces existen números enteros únicos $q$ y $r$ tales que
    $a = bq + r$ con $0 \le r < b$.
    \src{1-1.9}
\end{theorem}

\begin{idea}
    $bq$ es el mayor múltiplo de $b$ que no supera a $a$, y $r$ es lo que falta.
\end{idea}

\begin{proof}
    \emph{Existencia.} Sea $S = \{\, a - bk : k \in \mathbb{Z},\ a - bk \ge 0 \,\}$.
    $S$ no es vacío: para $k = -|a|$ se tiene $a + b|a| \ge a + |a| \ge 0$.
    Por el principio de buena ordenación, $S$ tiene un mínimo $r = a - bq \ge 0$.
    Si fuera $r \ge b$, entonces $r - b = a - b(q+1)$ estaría en $S$ y sería menor
    que $r$, lo que es absurdo. Luego $0 \le r < b$.

    \emph{Unicidad.} Si $a = bq + r = bq' + r'$ con $0 \le r, r' < b$, entonces
    $b(q - q') = r' - r$ y $|r' - r| < b$. Un múltiplo de $b$ con valor absoluto
    menor que $b$ solo puede ser $0$ (Proposición~\ref{prop:div-orden}), así que
    $r = r'$, y por la Proposición~\ref{prop:cancelacion}, $q = q'$.
\end{proof}

\begin{corollary}[Algoritmo de la División]\label{cor:division}
    Dados dos enteros $a$ y $b$ con $b \ne 0$ entonces existen $q$ y $r$ únicos tales que
    $a = bq + r$ donde $0 \le r < |b|$.
    \src{1-1.10}
\end{corollary}

\begin{idea}
    Aplicar el \nameref{teo:division} para $|b|$.
\end{idea}

\begin{definition}[Operador $\MOD$]\label{def:mod}
    Sean $a, b \in \mathbb{Z}$ con $b \neq 0$. Definimos $a \MOD b$ como el resto $r$
    que da el Corolario~\fullref{cor:division} al dividir $a$ entre $b$; es decir,
    el único $r$ con $0 \le r < |b|$ tal que $a = bq + r$ para algún $q \in \mathbb{Z}$.
    \src{1-1.14}
\end{definition}

\begin{proposition}[Propiedades del módulo]\label{prop:mod}
    Para cualesquiera $a, b, c, d, m \in \mathbb{Z}$ con $m \ne 0$.
    Si $a \MOD m = c \MOD m$ y $b \MOD m = d \MOD m$, entonces:
    \begin{enumerate}[label= (\arabic*), ref=\theproposition(\arabic*)]
        \item\label{prop:mod-suma} $(a+b) \MOD m = (c+d) \MOD m$.
        \item\label{prop:mod-producto} $(ab) \MOD m = (cd) \MOD m$.
        \src{1-1.16}
    \end{enumerate}
\end{proposition}

\begin{exercise}
    Para ningún número natural $n$ es $n^2 + 1$ divisible por 7.
    \src{1-1.18}
\end{exercise}

\begin{proof}
    $n^2 + 1$ es divisible entre 7 si y solo si $n^2 = 7k + 6$ para algún entero $k$, es decir,
    si y solo si $(n^2) \MOD 7 = 6$.
    Basta, por tanto, con ver que $n^2 \MOD 7 \neq 6$ para todo $n$.
    Observamos que:
    \[ n \MOD 7 = (n \MOD 7) \MOD 7 \]
    utilizando la Proposición~\ref{prop:mod-producto},
    \[ n^2 \MOD 7 = {(n \MOD 7)}^2 \MOD 7 \]
    ya que $0 \le n \MOD 7 < 7$, será suficiente comprobar los 7 valores posibles de $n \MOD 7$.
    \begin{align*}
        0^2 \MOD 7 &= 0, & 1^2 \MOD 7 &= 1, & 2^2 \MOD 7 &= 4, & 3^2 \MOD 7 &= 2, \\
        4^2 \MOD 7 &= 2, & 5^2 \MOD 7 &= 4, & 6^2 \MOD 7 &= 1.
    \end{align*}
    Ningún resultado es 6, luego $n^2 \MOD 7 \neq 6$ para todo $n$.
\end{proof}

\begin{definition}[Máximo Común Divisor]\label{def:mcd}
    Sean $a, b \in \mathbb{Z}$. Un \emph{divisor común} de $a$ y $b$ es un entero
    $d \neq 0$ tal que $d \mid a$ y $d \mid b$.

    El \emph{máximo común divisor} de $a$ y $b$, denotado $\mcd(a,b)$, es el divisor
    común $d > 0$ al que divide cualquier otro divisor común de $a$ y $b$.
    Por convenio, $\mcd(0,0) = 0$.

    Si $\mcd(a,b) = 1$, decimos que $a$ y $b$ son \emph{primos entre sí} o \emph{coprimos}.
    \src{1-1.19}
\end{definition}

\begin{theorem}[Identidad de Bézout]\label{teo:bezout}
    Sean $a$ y $b$ enteros distintos de $0$. Entonces existe un único $d$ máximo
    común divisor de $a$ y $b$. Además $d$ es el entero positivo más pequeño que
    puede expresarse en la forma $ax + by$ donde $x$ e $y$ son números enteros.
    \src{1-1.22}
\end{theorem}

\begin{proof}
    Probemos que existe $d > 0$ que divide a $a$ y $b$, y
    a su vez es divisible por cada común divisor de $a$ y $b$.

    Tomando cualquier divisor común $c$,
    por la Proposición~\ref{prop:div-comb-lineal}, $c \mid (ax + by)$
    para cualquier par de enteros $x$ e $y$.

    Ahora probemos la existencia del menor elemento que tenga la forma $ax + by$.
    Sea $S$ el conjunto de números $z$ que cumplen la propiedad $z = ax + by$ con $z > 0$.
    $S$ es un conjunto no vacío dado que, para $a > 0$ tomando el par $x = 1$ y $y = 0$
    resultaría en $z = a$, y para $a < 0$ tomando el par $x = -1$ y $y = 0$ resultaría en $z = -a$.
    Debido a que no es el conjunto vacío y está acotado inferiormente, por el principio
    de la buena ordenación existe un elemento mínimo de $S$, lo denotaremos por $d = ax_0 + by_0$.

    Utilizando el \nameref{cor:division}, $a = dq + r$, donde $0 \le r < d$.
    Procedamos por reducción al absurdo, supongamos $r > 0$.
    Podemos reescribir $r$ de la siguiente manera:
    \begin{align*}
        r &= a - dq = a - (ax_0 + by_0)q \\
          &= a(1 - x_0q) + b(-y_0q)
    \end{align*}
    y como $(1 - x_0q)$ y $(-y_0q)$ son enteros, $r \in S$. Pero $r < d$ llegando a una contradicción
    ya que $d$ es el menor elemento de $S$, por lo tanto $r = 0$, siendo así $d$ un divisor de $a$.

    Se demuestra la afirmación $d \mid b$ de forma análoga.
    Como $d$ es de la forma $ax + by$ todo divisor común a $a$ y $b$ divide a $d$.
    Por tanto, $d$ es un máximo común divisor de $a$ y $b$, lo que demuestra su existencia.

    Sea $d'$ otro máximo común divisor de $a$ y $b$.
    Como ambos son máximo común divisores, $d, d' > 0$.
    Como $d'$ es divisor común y $d$ es máximo común divisor, $d' \mid d$.
    Y recíprocamente $d \mid d'$.
    Cumplen las condiciones necesarias para aplicar la Proposición~\ref{prop:div-orden}: $d \le d'$ y $d' \le d$.
    Por lo tanto $d = d'$.

    Así, el máximo común divisor es único y coincide con $d$,
    que por construcción es el menor entero positivo de la forma $ax + by$.
\end{proof}


% Faltan teoremas por apuntar, se puede simplemente escribir la idea sin
% demostracion son teoremas menores



\subsection{Números Primos y Teorema Fundamental de la Aritmética}

\begin{lemma}[Lema de Euclides]\label{lem:euclides}
    Sean $a$, $b$ y $c$ números enteros. Supongamos que $a$ y $c$ son primos entre
    sí y que $c \mid ab$. Entonces $c \mid b$.
    \src{1-2.6}
\end{lemma}

\begin{proof}
    Si $a = 0$, entonces $1 = \mcd(a, c) = \mcd(0, c) = |c|$,
    luego $c = \pm1$ y $c \mid b$.

    Supongamos ahora $a \ne 0$. Como $c \mid ab$, se tiene $c \ne 0$.
    Por tanto, podemos aplicar el Teorema~\fullref{teo:bezout}: existen $s$ y $t$ enteros
    tales que $as + ct = 1$. Multiplicando ambos miembros de la igualdad por $b$ obtenemos:
    \[ abs + cbt = b. \]

    Como $c \mid ab$, por definición de divisibilidad existe un entero $k$ tal que
    $ab = ck$. Sustituyendo $ab$ en la igualdad anterior:
    \begin{align*}
        cks + cbt &= b \\
        c(ks + bt) &= b.
    \end{align*}
    Dado que $ks + bt$ es un número entero, $c$ divide a $b$.
\end{proof}